\begin{table}[!htbp]\centering\small
\caption{Implied contribution of immigration to national rent-index growth, 2015--2024, under explicit assumptions}\label{tab:magnitude}
\begin{adjustbox}{max width=\linewidth}\begin{tabular}{lccc}\toprule
Elasticity used & $\beta$ & Robust 95\% set for $\beta$ & Implied share of IPVA growth (\%)\\\midrule
Central specification, 2012--2024 & 0.13 & [$-$0.08, 0.35] & 4.1 \; [0, 11.0]\\
Central specification, arrival period 2015--2024 & 0.24 & [0.07, 0.59] & 7.5 \; [2.2, 18.6]\\
Central, renter weights, 2015--2024 & 0.20 & [0.06, 0.47] & 6.2 \; [1.9, 14.8]\\
Central + province $\times$ size FE, 2015--2024 & 0.13 & [$-$0.06, 0.52] & 4.2 \; [0, 16.4]\\
Original covariates, 2012--2024 & 0.30 & [0.09, 0.58] & 9.6 \; [2.8, 18.3]\\
Colombia--Peru episode (DiD-IV) & 0.64 & [0.04, 1.80] & 20.1 \; [1.3, 56.8]\\
\bottomrule\end{tabular}\end{adjustbox}
\begin{minipage}{0.97\linewidth}\footnotesize\vspace{4pt}
\textit{Notes}: implied share $=\beta\times M/G$, with $M=6.5$\,\% the cumulative net foreign-born inflow relative to population (2016--2024, centred timing) and $G=20.6$ log points the growth of the national IPVA between 2015 and 2024. The bracket translates the robust set into shares, truncating negative values at zero (the AR set with province clusters; wild-bootstrap set for the episode). The calculation assumes (A1) that the national elasticity equals the within-province elasticity, (A2) that level effects are permanent, (A3) that the 2015--2024 origin mix has the same elasticity as the identifying variation and (A4) no general-equilibrium effects; none of these is identified by the design (section~\ref{sec:magnitude}).
\end{minipage}\end{table}